% SOURCE_AUTHORITY: sga5_fr_workpass.tex, lines 7696--7735 (LF-normalized unit).
% SOURCE_SHA256: 791F4EFFC5E02832D5D77ED03518C8156D6F07E4C8238B03545DB93D883FBB28
\subsubsection*{5.2. Estrutura de \(J\)-adn\((\mathcal C)\) e AR-\(J\)-adn\((\mathcal C)\)}

Neste número, designaremos por \(\mathcal E_{\mathcal C}\), ou por \(\mathcal E\), se não houver risco de confusão, a subcategoria plena de \(\underline{\Hom}(\mathbb N^\circ,\mathcal C)\) cujos objetos são os sistemas projetivos AR-\(J\)-ádicos noetherianos.

\begin{statement}{Proposição 5.2.1}
A categoria \(\mathcal E_{\mathcal C}\) é estável por núcleos e conúcleos em \(\Hom(\mathbb N^\circ,\mathcal C)\). Em outras palavras, \(\mathcal E_{\mathcal C}\) é uma categoria abeliana e o funtor de inclusão
\[
\mathcal E_{\mathcal C}\longrightarrow \underline{\Hom}(\mathbb N^\circ,\mathcal C)
\]
É exato.
\end{statement}

\begin{prooftext}
Seja \(u:X\to Y\) um morfismo de \(\mathcal E_{\mathcal C}\). Veremos sucessivamente que o conúcleo, a imagem e o núcleo de \(u\) são AR-\(J\)-ádicos noetherianos. (Naturalmente, a demonstração de que a imagem de \(u\) é AR-\(J\)-ádica noetheriana é apenas uma etapa técnica intermediária).

a) \(\Coker(u)\) é AR-\(J\)-ádico noetheriano. Como \(X\) verifica (MLAR) e \(Y\) é AR-\(J\)-ádico, \(\Coker(u)\) é AR-\(J\)-ádico (3.2.4 (i)).

b) \(\Im(u)\) é AR-\(J\)-ádica noetheriana. Graças à sequência exata
\[
0\longrightarrow \Im(u)\longrightarrow Y\longrightarrow \Coker(u)\longrightarrow 0,
\]
A afirmação se deduzirá do lema 5.2.1.1 que aparece a seguir.
\end{prooftext}

\begin{statement}{Lema 5.2.1.1}
Seja
\[
0\longrightarrow L\longrightarrow M\longrightarrow N\longrightarrow0
\]
uma sequência exata na categoria \(\Hom(\mathbb N^\circ,\mathcal C)\). Se \(M\) e \(N\) são AR-\(J\)-ádicos noetherianos, \(L\) é AR-\(J\)-ádico noetheriano.
\end{statement}

\begin{prooftext}
Evidentemente, basta ver que é AR-\(J\)-ádico. Reduzimo-nos, da mesma forma que na demonstração de (3.2.4), ao caso em que \(M\) e \(N\) são estritos, e depois até mesmo \(J\)-ádicos noetherianos. Neste caso, de 3.1.3 (ii) deduz-se que \(L\) é estrito. O sistema projetivo \(L\) é estrito, está adaptado a \(J\), e seu graduado, como subobjeto graduado de \(\grs(M)\), é noetheriano: \(\grs(M)\), quociente de \(\gr_J(A)\otimes_A M_0\), é noetheriano (5.1.6). A afirmação deduz-se então do lema de Artin-Rees (4.2.6).
\end{prooftext}

c) \(\Ker(u)\) é AR-\(J\)-ádico noetheriano. Basta aplicar (5.2.1.1) à sequência exata evidente que se encontra abaixo:
\[
0\longrightarrow \Ker(u)\longrightarrow X\longrightarrow \Im(u)\longrightarrow0.
\]
